%=====================================================================
\section{Zone d'étude et données}

\subsection{La sylvo-écorégion G23 « Morvan et Autunois »}

La zone d'étude est celle du guide méthodologique de l'Irstea~\cite{ose2015guide} : la sylvo-écorégion
G23~\cite{ign2011ser}, massif de moyenne montagne (environ 200 à 900~m) de 3\,011~km² à cheval sur la Nièvre, l'Yonne, la
Côte-d'Or et la Saône-et-Loire. La forêt y couvre environ la moitié du territoire (52~\% d'après la BD Forêt). Le massif est connu pour ses
plantations de résineux, surtout de douglas, qui sont exploitées par coupes rases et remplacent progressivement la
forêt feuillue (chênaie, hêtraie). La BD Forêt V2~\cite{ign2024bdforet}, lue par le service WFS de la
Géoplateforme, donne 13\,494 polygones de forêt. En ne gardant que les forêts à couvert arboré et les peupleraies
(la classe FF0, \enquote{forêt fermée sans couvert}, correspond à des coupes anciennes), on obtient environ
157\,800~ha : 89\,750~ha de feuillus, 49\,400~ha de résineux et 18\,600~ha de forêts mixtes.

\subsection{Données}

\begin{table}[htbp]
\centering\small
\caption{Données utilisées sur le Morvan (juin--novembre 2026).}
\label{tab-data}
\begin{tabular}{p{2.7cm}p{5.4cm}p{6.7cm}}
\toprule
Capteur & Produits & Détails \\
\midrule
NISAR (L, 1,26~GHz) & L2 GCOV ($\gamma^0$ HH, HV, RTC) et GUNW (cohérence HH à 20~m), produits par
ISCE3~\cite{isce3} & 2 géométries couvrant la fenêtre radar : A173~F25 (montante) et D66~F65 (descendante). Paires
de 12 à 24~jours. \\
Sentinel-1 (C, 5,4~GHz) & IW SLC, VV+VH, traités avec SNAP~\cite{snap} & Orbites A59 (IW3) et D37 (IW2), 3 bursts
chacune. 15 dates à 12~jours par orbite, du 01/06 au 18/11/2026, presque toutes de S1D. \\
Sentinel-2~\cite{drusch2012sentinel2} & L2A, NDVI et SWIR & Référence : méthode du guide sur toute la G23 (tuiles
31TEM, 31TEN, 31TFM, 31TFN) et photo-interprétation. \\
IGN & Contour G23, BD Forêt V2 & Masque forêt et groupes (feuillus, résineux, mixtes). \\
\bottomrule
\end{tabular}
\end{table}

Comme dans les Vosges, les volumes téléchargés sont réduits : les produits NISAR sont lus à distance par blocs
HDF5, et seuls les fichiers de la sous-fauchée utile sont extraits des archives Sentinel-1.

%=====================================================================
\section{Méthode}

\subsection{La méthode du guide}

Le guide~\cite{ose2015guide} détecte les coupes rases en six étapes : NDVI de deux images d'été, différence
$\Delta$NDVI, masques zone d'étude et forêt, seuillage statistique dans la forêt, tamisage à 1~ha et évaluation par
matrice de confusion en surface. Le seuillage utilise la moyenne $m$ et l'écart-type $\sigma$ de $\Delta$NDVI
dans la forêt, avec trois degrés de certitude :
\begin{center}\small
\begin{tabular}{ll}
\toprule
Degré & Intervalle \\
\midrule
1 (faible) & $[m-2\sigma, m-\sigma[$ \\
2 (moyen) & $[m-3\sigma, m-2\sigma[$ \\
3 (fort) & $< m-3\sigma$ \\
\bottomrule
\end{tabular}
\end{center}
Ses spécifications sont une unité minimale de 1~ha et des erreurs d'omission et de commission $\leq 10$~\% en
surface. Le module \texttt{nisar\_s1.guide\_cr} implémente ces étapes pour n'importe quelle image de changement
($\Delta$NDVI, $\Delta\gamma^0$, baisse de cohérence), avec le sens du changement en paramètre.

\subsection{Référence optique}

\paragraph{Guide appliqué à Sentinel-2.} Pour chaque tuile, le NDVI est la médiane des 3 dates les plus claires de
deux fenêtres : \enquote{avant} (15/06--08/07/2026) et \enquote{après} (15/09--16/10/2026). Les seuils du guide
sont calculés dans la forêt de G23.

\paragraph{Référence qualifiée.} Le guide ne retient que les parcelles \enquote{passées d'un état boisé à un état
sol nu}. La référence est donc : degré~3, NDVI avant $\geq 0{,}70$, NDVI après $\leq 0{,}65$, puis tamisage à
1~ha.

\paragraph{Photo-interprétation.} Comme le recommande le guide (section 3.9.1), la référence est contrôlée
visuellement : couleurs vraies avant et après, et fausses couleurs SWIR / PIR / rouge après, où le sol nu et les
rémanents apparaissent rose-magenta. Chaque objet reçoit un verdict : coupe certaine, probable ou sans sol nu.

\paragraph{Test \enquote{forêt avant}.} Une coupe rase est la perte d'un peuplement fermé. La réflectance SWIR
(B11) de juillet dans l'objet est comparée à celle d'un anneau de forêt de 100 à 400~m : l'objet était un
peuplement fermé si sa B11 ne dépasse pas le 90\textsuperscript{e} centile de l'anneau.

\subsection{Fenêtre radar}

Le traitement Sentinel-1 se limite à 3 bursts d'une seule sous-fauchée (mémoire vive). Les fenêtres de
$20\times20$~km ont été classées selon la surface de coupes de référence, puis contrôlées dans les catalogues ASF :
une orbite Sentinel-1 montante et une descendante, chacune dans une seule sous-fauchée, et au moins une géométrie
NISAR montante et une descendante. La fenêtre retenue (nord-est de G23) permet, contrairement aux Vosges, la
comparaison \textbf{à géométrie égale dans les deux directions} : A173 (L) contre A59 (C) en montante, D66 (L)
contre D37 (C) en descendante.

\subsection{Traitements radar}

\begin{itemize}
\item \textbf{NISAR} : $\gamma^0$ HH et HV des GCOV (RTC~\cite{small2011flattening}) et cohérence HH des GUNW,
  utilisés tels que livrés par ISCE3 et rééchantillonnés sur la grille Lambert 93 de 20~m.
\item \textbf{Sentinel-1} : mêmes graphes SNAP que pour les Vosges. Pour $\gamma^0$ : calibration,
  Terrain-Flattening et géocodage. Pour la cohérence : fenêtre $10\times3$, phases plate et topographique
  retirées. Les 3 bursts sont raccordés par TOPSAR-Deburst, sans ESD.
\end{itemize}

\subsection{Variables de changement}

Les acquisitions sont classées en \enquote{avant} (01/06--12/07), \enquote{après} (10/09--20/10) et
\enquote{entre} ces deux fenêtres, qui encadrent les fenêtres optiques. Les dates de coupe ne sont pas connues
une à une : les variables sont donc définies pour être indépendantes de la date exacte.
\begin{itemize}
\item $\Delta\gamma^0$ : $\gamma^0$ moyen après $-$ $\gamma^0$ moyen avant (dB, moyennes en puissance linéaire) ;
\item $\gamma^0$ après (dB) ;
\item cohérence minimale sur les paires \enquote{entre} : la paire qui encadre la coupe en fait partie, quelle que
  soit la date de la coupe ;
\item chute de cohérence (cohérence avant $-$ cohérence minimale) et cohérence après.
\end{itemize}

\subsection{Évaluation}

\paragraph{Classes.} La classe \enquote{coupe} est formée des pixels forestiers des coupes de référence. La classe
\enquote{forêt intacte} est la forêt BD Forêt dont $\Delta$NDVI reste dans $m \pm \sigma$ et située à plus de 40~m
de toute détection optique. Les objets écartés par la photo-interprétation et une zone tampon de 40~m sont exclus
des deux classes.

\paragraph{Séparabilité.} Pour chaque variable : AUC (statistique de Mann-Whitney~\cite{hanley1982auc}, exprimée
comme $\max(\mathrm{AUC}, 1-\mathrm{AUC})$), d de Cohen~\cite{cohen1988power} et distance de
Jeffries-Matusita~\cite{swain1978remote}. Les intervalles de confiance à 95~\% sont obtenus par un bootstrap en
grappes~\cite{kunsch1989bootstrap} : on rééchantillonne les coupes entières (et non les pixels, qui sont
corrélés) et des blocs de $5\times5$ pixels de forêt. Avec 7 coupes, ces intervalles sont larges.

\paragraph{Contrôles.} Avant la coupe, coupes et forêt doivent être indiscernables. On le vérifie par le
changement entre deux dates \enquote{avant} et par le niveau de $\gamma^0$ avant la coupe, qui révèle un éventuel
effet du type de peuplement.

\paragraph{Transposition du guide au radar.} Les degrés de certitude $m-k\sigma$ sont calculés sur chaque variable
radar dans la forêt de la fenêtre, puis tamisés à 1~ha. On calcule ensuite la matrice de confusion en surface du
guide contre la référence (omission, commission, kappa) et la part des coupes trouvées (au moins 25~\% de leurs
pixels détectés). On calcule aussi le taux de fausses alertes rapporté à la surface forestière : avec une
référence de quelques hectares, la commission relative du guide est très sensible à quelques hectares de fausses
alertes.

\paragraph{Audit des fausses alertes.} Les plus grandes détections radar hors de toute détection optique sont
photo-interprétées sur une image Sentinel-2 tardive (novembre) : une partie peut être des coupes réalisées après
la fenêtre optique de septembre.

\paragraph{Datation.} Pour chaque coupe de référence et chaque géométrie, la série de $\gamma^0$ en polarisation
croisée est comparée à son niveau d'avant juillet. La coupe est datée entre la dernière acquisition normale et la
première acquisition d'une baisse durable de plus de 2~dB (deux acquisitions consécutives).
